\documentclass[12pt,a4paper]{article}

% Paquetes básicos
\usepackage[utf8]{inputenc}   % Acentos y caracteres especiales
\usepackage[T1]{fontenc}
\usepackage[spanish]{babel}   % Idioma español
\usepackage{graphicx}         % Insertar imágenes
\usepackage{amsmath, amssymb} % Matemáticas
\usepackage{hyperref}         % Hipervínculos
\usepackage{csquotes}         % Citas
\usepackage{geometry}         % Márgenes
\geometry{margin=2.5cm}

% Datos del artículo
\title{Título del Artículo}
\author{Nombre del Autor \\[0.5em]
	\small Universidad / Institución \\
	\small \texttt{correo@ejemplo.com}}
\date{\today}

\begin{document}
	
	% Portada
	\maketitle
	
	% Resumen
	\begin{abstract}
		Este es un breve resumen del artículo. Aquí se expone de manera concisa el objetivo, la metodología y los principales hallazgos o aportes del trabajo.
	\end{abstract}
	
	% Palabras clave
	\textbf{Palabras clave:} palabra1, palabra2, palabra3
	
	\section{Modelado de Tópicos de todas las palabras del corpus}
	Paper de modelado de tópicos con BERTopic https://arxiv.org/abs/2203.05794
	
	Github
	https://github.com/MaartenGr/BERTopic\\
	
	
	- modelado por año
	- modelado por mes
	
	
	\section{Introducción}
	En esta sección se presenta el contexto, la motivación y los objetivos del trabajo.  
	Se puede citar con el comando \cite{autor2020}.
	
	\section{Marco Teórico}
	Aquí se revisan los conceptos, teorías o investigaciones previas que fundamentan el artículo.
	
	\section{Metodología}
	Se detalla el enfoque, los métodos, datos o experimentos utilizados.
	
	\section{Resultados}
	Presentación y análisis de los resultados obtenidos.  
	Puedes incluir una figura (Figura %\ref{fig:ejemplo}):
	
	
	
	
	
	\section{Discusión}
	Interpretación de los resultados, comparación con otros estudios y relevancia.
	
	\section{Conclusiones}
	Se resumen los hallazgos principales, limitaciones y posibles líneas de investigación futura.
	
	% Bibliografía
	\begin{thebibliography}{9}
		\bibitem{autor2020}
		Apellido, N. (2020). \textit{Título del libro o artículo}. Editorial o revista.
		
		\bibitem{otro2021}
		Apellido, M. (2021). \textit{Otro recurso citado}. Editorial o revista.
	\end{thebibliography}
	
\end{document}
